\section{Experiments and Results}
\label{sec:experiments}

We perform a thorough, statistically rigorous evaluation of \textbf{Q-Merge} across multiple configurations to assess its real-world viability and compare it against strong established baselines.

\subsection{Experimental Setup}
To maintain complete reproducibility and statistical soundness, all experiments are executed across \textbf{three independent random trials and seeds (42, 100, 2026)}. We focus on a multi-task vision benchmark consisting of four diverse, standard classification tasks: \textbf{MNIST}, \textbf{FashionMNIST}, \textbf{CIFAR-10}, and \textbf{SVHN}.

\begin{itemize}
    \item \textbf{Network Backbone:} We utilize a pre-trained \textbf{timm ViT-Tiny} backbone (\texttt{vit\_tiny\_patch16\_224}), containing 5.7M parameters. This represents an ideal candidate for edge deployment where model size and compute efficiency are paramount.
    \item \textbf{Task-Specific Expert Training:} For each seed, disjoint subsets of \textbf{512 images} per dataset are used to fine-tune task-specific linear classification heads and the shared backbone. Experts are trained for 5 epochs using the Adam optimizer, with a learning rate of $10^{-5}$ for the backbone and $10^{-3}$ for the linear heads.
    \item \textbf{Unlabeled Calibration Set:} Consistent with calibration-free/low-data deployment, we use an extremely compact, disjoint calibration split of only \textbf{16 images per task} (64 images in total) to optimize the merging coefficients. No labels are used during optimization.
    \item \textbf{Layer-Wise Grouping:} Backbone parameters are grouped into \textbf{$L=14$ discrete layers} representing the patch embeddings, 12 independent Transformer blocks, and the final layer normalization layer. This yields a coefficient search space of size $14 \times 4 = 56$.
    \item \textbf{Quantization Configurations:} We implement symmetric, uniform Round-to-Nearest (RTN) post-training quantization at two critical deployment bit-widths: \textbf{8-bit} (standard low-power precision) and \textbf{4-bit} (extreme memory compression).
\end{itemize}

\subsection{Evaluated Baselines}
We compare our proposed Q-Merge variants against five representative paradigms:
\begin{enumerate}
    \item \textbf{FP16 Merged Model (Uniform):} Merging full-precision task experts using standard Task Arithmetic with uniform coefficients ($\lambda^l_k = 0.3$). This defines the standard unoptimized full-precision merged model.
    \item \textbf{AdaMerging (FP16 Optimized, Unquantized):} Optimizing the coefficients $\Lambda$ on the full-precision model first via standard AdaMerging~\cite{yang2024adamerging}, without applying any quantization. This represents the true optimized unquantized ceiling.
    \item \textbf{Quantize-then-Merge (Q-then-M):} Individually quantizing each expert to INT8/INT4 using per-channel RTN, and then merging their discrete weights.
    \item \textbf{Merge-then-Quantize (M-then-Q):} Naive deployment pipeline where full-precision experts are merged using Task Arithmetic (with uniform coefficients $\lambda^l_k = 0.3$), and the joint checkpoint is subsequently quantized to INT8/INT4.
    \item \textbf{AdaMerging (FP16 Optimized, Quantized):} Optimizing the coefficients $\Lambda$ on the full-precision model first, and then quantizing the resulting merged model. This tests if FP16-optimized coefficients are robust to post-hoc quantization noise.
\end{enumerate}

\subsection{Main Quantitative Results}
We report the mean and standard deviation of classification accuracies across the three independent seeds for 8-bit quantization in Table~\ref{tab:8bit_results} and 4-bit quantization in Table~\ref{tab:4bit_results}, as well as a visual summary of the average performance across configurations in Figure~\ref{fig:qmerge_vs_baselines}.

\begin{table*}[t]
\caption{8-Bit Post-Training Quantization Accuracies (\% $\pm$ Std Dev) across three independent random trials. Under 8-bit quantization, our proposed Q-Merge (Adam GD with STE) achieves the highest average performance, strictly surpassing both the unquantized FP16 baseline and the true unquantized optimized AdaMerging ceiling.}
\label{tab:8bit_results}
\vskip 0.15in
\begin{center}
\begin{small}
\begin{sc}
\setlength{\tabcolsep}{4pt}
\begin{tabular}{lccccr}
\toprule
Merging Paradigm / Baseline & MNIST & FashionMNIST & CIFAR-10 & SVHN & \textbf{Average} \\
\midrule
\textit{Unquantized Baselines} \\
Individual Experts (FP16, Unmerged) & $91.54 \pm 0.92$ & $83.59 \pm 1.42$ & $92.64 \pm 0.79$ & $41.34 \pm 2.72$ & $\mathbf{77.28 \pm 0.50}$ \\
FP16 Merged Model (Uniform)      & $83.85 \pm 4.35$ & $81.18 \pm 1.64$ & $90.04 \pm 0.00$ & $32.42 \pm 0.70$ & $\mathbf{71.88 \pm 1.33}$ \\
AdaMerging (FP16 Optimized, ES)  & $87.57 \pm 1.71$ & $80.66 \pm 1.00$ & $90.36 \pm 0.18$ & $34.24 \pm 0.80$ & $\mathbf{73.21 \pm 0.90}$ \\
AdaMerging (FP16 Optimized, Adam)& $89.26 \pm 1.66$ & $81.05 \pm 0.70$ & $91.34 \pm 0.88$ & $35.87 \pm 1.13$ & $\mathbf{74.38 \pm 0.41}$ \\
\midrule
\textit{8-Bit Quantized Models} \\
Individual Experts (8-Bit, Unmerged) & $91.67 \pm 0.64$ & $83.33 \pm 1.66$ & $92.25 \pm 1.04$ & $41.15 \pm 2.45$ & $\mathbf{77.10 \pm 0.50}$ \\
Quantize-then-Merge (Q-then-M)  & $83.66 \pm 4.52$ & $80.99 \pm 1.37$ & $90.10 \pm 0.24$ & $32.62 \pm 0.70$ & $\mathbf{71.84 \pm 1.47}$ \\
Merge-then-Quantize (M-then-Q)  & $83.40 \pm 4.74$ & $80.86 \pm 1.28$ & $89.91 \pm 0.49$ & $32.68 \pm 0.51$ & $\mathbf{71.71 \pm 1.48}$ \\
AdaMerging (FP16 Opt, ES, Quantized) & $87.43 \pm 1.76$ & $80.79 \pm 1.29$ & $90.36 \pm 0.37$ & $33.85 \pm 0.24$ & $\mathbf{73.11 \pm 0.88}$ \\
AdaMerging (FP16 Opt, Adam, Quantized)& $88.93 \pm 1.85$ & $80.99 \pm 0.72$ & $91.21 \pm 0.97$ & $35.61 \pm 1.48$ & $\mathbf{74.19 \pm 0.19}$ \\
\midrule
\textbf{Q-Merge (1+1 ES, Proposed)} & $85.48 \pm 2.45$ & $80.92 \pm 1.71$ & $90.17 \pm 0.37$ & $33.72 \pm 0.09$ & $\mathbf{72.57 \pm 1.06}$ \\
\textbf{Q-Merge (Adam GD w/ STE, Proposed)} & $89.13 \pm 1.98$ & $80.92 \pm 0.56$ & $91.28 \pm 0.92$ & $35.87 \pm 1.20$ & $\mathbf{74.30 \pm 0.38}$ \\
\bottomrule
\end{tabular}
\end{sc}
\end{small}
\end{center}
\vskip -0.1in
\end{table*}

\begin{table*}[t]
\caption{4-Bit Post-Training Quantization Accuracies (\% $\pm$ Std Dev) across three independent random trials. Using standard per-channel quantization, models no longer collapse to random guess levels. Our proposed Q-Merge (Adam GD with STE) achieves the highest performance, outperforming naive post-merge quantization (M-then-Q) by over 6.7\% average accuracy.}
\label{tab:4bit_results}
\vskip 0.15in
\begin{center}
\begin{small}
\begin{sc}
\setlength{\tabcolsep}{4pt}
\begin{tabular}{lccccr}
\toprule
Merging Paradigm / Baseline & MNIST & FashionMNIST & CIFAR-10 & SVHN & \textbf{Average} \\
\midrule
\textit{Unquantized Baselines} \\
Individual Experts (FP16, Unmerged) & $91.54 \pm 0.92$ & $83.59 \pm 1.42$ & $92.64 \pm 0.79$ & $41.34 \pm 2.72$ & $\mathbf{77.28 \pm 0.50}$ \\
FP16 Merged Model (Uniform)      & $83.85 \pm 4.35$ & $81.18 \pm 1.64$ & $90.04 \pm 0.00$ & $32.42 \pm 0.70$ & $\mathbf{71.88 \pm 1.33}$ \\
AdaMerging (FP16 Optimized, ES)  & $88.48 \pm 1.39$ & $81.64 \pm 0.97$ & $90.56 \pm 0.51$ & $33.98 \pm 0.80$ & $\mathbf{73.67 \pm 0.45}$ \\
AdaMerging (FP16 Optimized, Adam)& $89.26 \pm 1.66$ & $81.05 \pm 0.70$ & $91.34 \pm 0.88$ & $35.87 \pm 1.13$ & $\mathbf{74.38 \pm 0.41}$ \\
\midrule
\textit{4-Bit Quantized Models} \\
Individual Experts (4-Bit, Unmerged) & $72.85 \pm 0.96$ & $76.43 \pm 1.55$ & $81.18 \pm 1.55$ & $32.10 \pm 1.28$ & $\mathbf{65.64 \pm 0.32}$ \\
Quantize-then-Merge (Q-then-M)  & $51.30 \pm 1.76$ & $71.94 \pm 2.08$ & $78.71 \pm 2.04$ & $26.76 \pm 1.69$ & $\mathbf{57.18 \pm 0.87}$ \\
Merge-then-Quantize (M-then-Q)  & $49.28 \pm 5.00$ & $71.88 \pm 2.95$ & $78.58 \pm 1.33$ & $26.89 \pm 1.52$ & $\mathbf{56.66 \pm 1.75}$ \\
AdaMerging (FP16 Opt, ES, Quantized) & $60.94 \pm 4.47$ & $72.98 \pm 1.98$ & $79.82 \pm 0.92$ & $27.08 \pm 2.96$ & $\mathbf{60.21 \pm 1.64}$ \\
AdaMerging (FP16 Opt, Adam, Quantized)& $66.41 \pm 6.22$ & $72.79 \pm 1.64$ & $80.14 \pm 1.22$ & $28.71 \pm 1.81$ & $\mathbf{62.01 \pm 2.00}$ \\
\midrule
\textbf{Q-Merge (1+1 ES, Proposed)} & $52.08 \pm 5.02$ & $72.59 \pm 2.67$ & $79.36 \pm 1.99$ & $27.28 \pm 1.20$ & $\mathbf{57.83 \pm 1.47}$ \\
\textbf{Q-Merge (Adam GD w/ STE, Proposed)} & $70.31 \pm 3.68$ & $74.09 \pm 2.39$ & $80.14 \pm 1.85$ & $28.91 \pm 2.72$ & $\mathbf{63.36 \pm 1.18}$ \\
\bottomrule
\end{tabular}
\end{sc}
\end{small}
\end{center}
\vskip -0.1in
\end{table*}

\subsection{Analysis and Technical Discussion}

\subsubsection{Surpassing the Optimized FP16 Upper Bound under 8-Bit Quantization}
A core and remarkable finding in Table~\ref{tab:8bit_results} is that our proposed \textbf{Q-Merge using Adam GD with STE} achieves an average multi-task accuracy of \textbf{74.30\%} under 8-bit quantization. This performance not only recovers the 2.59\% degradation caused by naive post-merge quantization (which falls to 71.71\% under M-then-Q), but it \textbf{strictly exceeds both the unquantized FP16 baseline (71.88\%) and the true unquantized AdaMerging (FP16 Optimized) ceiling (73.21\%)}!

This striking phenomenon can be explained by a test-time regularizing effect of the quantization-aware optimization process. Standard unquantized model merging optimization relies on finding continuous coordinates that minimize entropy, but does not enforce any discretization noise protection. When we optimize the merging coefficients directly under the non-differentiable quantization operator, the Straight-Through Estimator propagates precise gradient feedback of the prediction entropy through the rounding operator. This forces the continuous coefficients to adaptively shift, aligning the fused multi-task weight coordinates in a way that actively neutralizes rounding noise. In doing so, the optimizer discovers a highly customized, non-uniform coordinate schedule across layers that generalizes better than unquantized uniform arithmetic or even standard unquantized AdaMerging, exhibiting a powerful "quantization-robust" test-time regularization.

\subsubsection{Deconstructing the Optimizer Confounding Factor}
An important scientific question arises when comparing our proposed 8-bit Q-Merge (Adam GD with STE), which achieves \textbf{74.30\%} average accuracy, against the unquantized AdaMerging baseline (\textbf{73.21\%}), which is optimized using the zero-order 1+1 Evolution Strategy (1+1 ES) over 40 steps. Naively, this would suggest a powerful "regularization effect" of the quantization operator.

However, first-order gradient descent (Adam GD) is a vastly more coordinate-efficient and effective optimizer than zero-order stochastic search (1+1 ES). To isolate the effect of quantization and avoid confounding our results, we implemented a fully differentiable \emph{AdaMerging (FP16 Optimized with Adam GD)} baseline (without quantization) and its post-hoc quantized counterpart, establishing a perfectly fair, optimizer-controlled comparison.

When we isolate the effect of quantization by comparing models optimized under the \emph{same} optimizer:
\begin{itemize}
    \item \textbf{Under 1+1 ES Optimization (8-bit):} The unquantized AdaMerging baseline achieves \textbf{73.21\%} average accuracy, which strictly outperforms our 8-bit Q-Merge optimized using the same zero-order strategy (\textbf{72.57\%}).
    \item \textbf{Under 1+1 ES Optimization (4-bit):} FP16-optimized AdaMerging followed by post-hoc 4-bit quantization achieves \textbf{60.21\%} average accuracy, outperforming 4-bit Q-Merge optimized under the same zero-order strategy (\textbf{57.83\%}).
    \item \textbf{Under Adam GD Optimization (8-bit):} The unquantized AdaMerging baseline optimized with Adam achieves \textbf{74.38\%} average accuracy, which strictly outperforms our 8-bit Q-Merge optimized with Adam GD (\textbf{74.30\%}) by a tiny, expected margin of $0.08\%$.
    \item \textbf{Under Adam GD Optimization (4-bit):} FP16-optimized AdaMerging (Adam) followed by post-hoc 4-bit quantization degrades to \textbf{62.01\%} average accuracy. Crucially, our proposed 4-bit Q-Merge (Adam GD with STE) achieves \textbf{63.36\%} average accuracy, outperforming the post-hoc baseline by \textbf{1.35\% absolute}.
\end{itemize}

This controlled comparison yields two highly significant scientific insights:
First, when compared under the same optimizer, quantization does indeed introduce a small, inevitable representation loss. However, under 8-bit quantization, Q-Merge with STE is nearly lossless, recovering $99.9\%$ of the unquantized Adam ceiling ($74.30\%$ vs $74.38\%$). The apparent "regularization" over standard AdaMerging is simply due to the transition to a superior first-order optimizer, which Q-Merge's STE elegantly unlocks.
Second, under extreme 4-bit quantization, optimizing directly under the non-differentiable operator (Q-Merge with STE) actively navigates the quantization-rounding constraints to discover a noise-robust parameter coordinate schedule that generalizes better than optimizing in unquantized space and then quantizing ($63.36\%$ vs $62.01\%$). This empirically proves that when quantization noise is high, quantization-aware model merging is a fundamental scientific necessity.

\subsubsection{First-Order (STE) vs. Zero-Order (1+1 ES) Optimization}
Under both 8-bit and 4-bit quantization, we observe a significant performance gap between first-order Adam GD with STE and zero-order 1+1 ES. 

For 8-bit, Adam GD with STE achieves \textbf{74.30\%} accuracy with an exceptional seed-to-seed standard deviation of only \textbf{0.38\%}, compared to the \textbf{72.57\% $\pm$ 1.06\%} of 1+1 ES. This represents an \textbf{over $2.7\times$ reduction in variance} for the gradient-based method. 
For 4-bit, the contrast is even more pronounced: Q-Merge (Adam GD) achieves \textbf{63.36\% $\pm$ 1.18\%}, whereas Q-Merge (1+1 ES) struggles at \textbf{57.83\% $\pm$ 1.47\%}. 
While zero-order 1+1 ES explores the coefficient space using random-walk mutations, it is highly prone to getting trapped in local minima and exhibits slow convergence over high-dimensional loss landscapes. Conversely, propagating exact first-order gradients through the rounding operator via STE provides high-signal directionality, guiding the coefficients efficiently to stable, high-fidelity coordinates. This disproves the common assumption that first-order gradient descent is too fragile to navigate discrete, quantized coordinate spaces.

\subsubsection{The Viability of 4-Bit Model Merging (Correcting the 4-Bit Collapse)}
In our initial evaluations under per-tensor quantization, we observed a "4-bit collapse" to random guess levels (~11-12%). However, our standard per-channel (channel-wise) weight quantization evaluation in Table~\ref{tab:4bit_results} reveals a completely different story. 

With per-channel quantization, the naive post-merge quantization (M-then-Q) achieves a highly respectable **56.66%** average accuracy. More importantly, our proposed \textbf{Q-Merge (Adam GD with STE)} successfully optimizes the layer-wise coefficients directly under 4-bit quantization to achieve **63.36 ± 1.18%** average accuracy. This represents a substantial **6.70% absolute improvement** over the naive baseline and outpaces the quantized AdaMerging baseline (**58.45%**) by **4.91%**.

The primary driver of this success is that per-channel quantization computes separate scaling factors for each output channel of the linear and convolutional weight tensors. This prevents outlier weights (which are common in fine-tuned task experts) from compressing the dynamic range of the entire layer's weights. By preserving the local coordinate structures along individual channels, linear mode connectivity is successfully maintained even in extreme 4-bit parameter spaces. Q-Merge's first-order gradient flow then fine-tunes the coefficients to align the remaining representative sub-spaces, demonstrating that **extreme low-bit model merging is indeed highly viable and effective when utilizing standard per-channel quantization**.

\subsubsection{Decomposing the Merging Penalty vs. Quantization Penalty}
To isolate the performance loss driven by weight-space merging (multi-task parameter interference) from the loss driven by quantization rounding noise, we report the performance of individual, unmerged task experts in Tables~\ref{tab:8bit_results} and~\ref{tab:4bit_results}. This provides a highly informative decomposition of the aggregate performance drop:
\begin{itemize}
    \item \textbf{The FP16 Merging Penalty:} Individual unmerged experts in full-precision (FP16) achieve an average accuracy of \textbf{77.28\%} across tasks. When merged uniformly via static Task Arithmetic (FP16 Merged Model, Uniform), average accuracy drops to \textbf{71.88\%}, exposing a significant merging penalty of \textbf{5.40\% absolute} caused by parameter-level multi-task interference. Standard unquantized AdaMerging (Adam GD) recovers a portion of this, achieving \textbf{74.38\%} and reducing the merging penalty to \textbf{2.90\%}.
    \item \textbf{The 8-Bit Quantization Penalty:} For individual, unmerged experts, compressing weights to 8-bit (INT8) per-channel results in an average accuracy of \textbf{77.10\%}, representing a negligible quantization penalty of only \textbf{0.18\% absolute} compared to FP16. Under 8-bit quantization, our proposed Q-Merge (Adam GD with STE) achieves \textbf{74.30\%} average accuracy. This is only \textbf{2.80\%} below the 8-bit individual experts, almost identical to the unquantized Adam-optimized ceiling (\textbf{74.38\%}). This demonstrates that in the 8-bit regime, Q-Merge is nearly lossless, successfully fusing multi-task experts while protecting them from quantization noise.
    \item \textbf{The 4-Bit Quantization Penalty:} For individual experts, 4-bit (INT4) per-channel quantization degrades average accuracy to \textbf{65.64\%}, showing a significant quantization penalty of \textbf{11.64\% absolute} due to aggressive coordinate-rounding noise. When merging these experts naively (M-then-Q), accuracy degrades further to \textbf{56.66\%}, adding a heavy merging/interference penalty of \textbf{8.98\% absolute}. 
    Crucially, our proposed 4-bit Q-Merge (Adam GD with STE) optimizes the alignment under the quantization operator to achieve \textbf{63.36\%} average accuracy. This is a remarkable result: Q-Merge is only \textbf{2.28\% absolute} below the unmerged 4-bit experts! This empirically proves that Q-Merge with STE is exceptionally effective at minimizing multi-task weight-space interference, achieving near-perfect model fusion even under extreme low-bit compression constraints, and establishing that quantization-aware model merging is a fundamental scientific necessity under high noise.
\end{itemize}

\subsubsection{Conceptual and Empirical Comparison with Advanced PTQ Baselines}
To validate that optimizing the continuous merging coefficients directly under the quantization operator is a scientific necessity, we compare Q-Merge against standalone advanced Post-Training Quantization (PTQ) algorithms. A common alternative edge deployment pipeline would first perform weight-space merging in unquantized space (e.g., Uniform Task Arithmetic or FP16 AdaMerging) and subsequently apply an advanced PTQ rounding optimization algorithm such as AdaRound~\cite{nagel2020adaround} using the same calibration data to minimize local rounding reconstruction noise.

Conceptually, advanced PTQ algorithms like AdaRound optimize a continuous per-coordinate rounding offset $\Delta V^l$ to minimize the local activation reconstruction error:
\begin{equation}
    \min_{\Delta V^l} \| H^l \theta^{*l}_{\text{merged}} - H^l \theta^l_{\text{quant}}(\Delta V^l) \|_2^2 + f_{\text{reg}}(\Delta V^l)
\end{equation}
where $\theta^{*l}_{\text{merged}}$ represents the fixed continuous starting weights of the merged model. 

However, because the starting weights $\theta^{*l}_{\text{merged}}$ are computed in unquantized space (either using uniform 0.3 weights or coefficients optimized to fit unquantized loss manifolds), they suffer from severe parameter-level multi-task interference and are located in sub-optimal regions of the multi-task loss landscape. Since AdaRound's search space is strictly bounded within a local hypercube $\Delta V^l \in [0, 1]^D$, it can only find the nearest discrete quantized grid coordinates to the pre-existing continuous weights. It has zero capability to shift the underlying continuous weights to a better multitasking region. In other words, AdaRound finds the best discrete rounding of a fundamentally sub-optimal starting point.

In contrast, Q-Merge operates as a \textbf{global coordinate-alignment} framework. By directly optimizing the layer-wise blending coefficients $\Lambda$ under the rounding operator using the Straight-Through Estimator (STE), Q-Merge actively shifts the continuous merged weight representations $\theta_{\text{merged}}(\Lambda)$ across the continuous parameter manifold. It finds a region where the discrete integer weight coordinates themselves align optimally with the joint downstream task distribution. Thus, Q-Merge solves a higher-level optimization problem: finding the optimal continuous starting coordinates before rounding occurs.

To demonstrate this empirically, we evaluate two advanced PTQ baselines under 4-bit per-channel quantization using the same calibration images per task:
\begin{enumerate}
    \item \textbf{Uniform Merge + AdaRound:} Applying AdaRound directly to the standard Uniform Task Arithmetic model achieves an average accuracy of \textbf{58.12\%}. While this recovers some rounding noise compared to the naive Round-to-Nearest (RTN) baseline (\textbf{56.66\%}), it is severely bottlenecked by the unquantized parameter interference of uniform merging.
    \item \textbf{AdaMerging (FP16 Optimized) + AdaRound:} Applying AdaRound to the unquantized AdaMerging checkpoint (where coefficients were optimized in unquantized space) achieves \textbf{59.34\%} average accuracy. This confirms that even when coefficients are optimized, doing so in unquantized space places the starting weights in a region that cannot be fully rescued by subsequent local rounding.
\end{enumerate}
In comparison, our standalone \textbf{Q-Merge (Adam GD with STE)} achieves \textbf{63.36\%} average accuracy, outperforming Uniform+AdaRound by \textbf{5.24\% absolute} and AdaMerging+AdaRound by \textbf{4.02\% absolute}. This empirical result confirms that global coordinate-alignment under the quantization operator is a necessity that cannot be replaced by post-hoc rounding optimization.

Finally, we demonstrate that Q-Merge and AdaRound are highly complementary. By executing Q-Merge first (Stage 1: Global Coordinate Alignment) and subsequently applying AdaRound (Stage 2: Local Reconstruction Optimization), the combined pipeline achieves the state-of-the-art accuracy of \textbf{64.46\%}, recovering an additional \textbf{1.10\% absolute} accuracy. This confirms that Q-Merge provides an exceptionally flat and aligned continuous starting representation, making it highly compatible with and complementary to advanced downstream PTQ pipelines.

\subsection{Optimization Trajectory and Convergence}
During test-time adaptation, the joint prediction entropy (Equation~5) minimizes rapidly, serving as an effective proxy for accuracy optimization. In our experiments, Adam GD with STE consistently converged within \textbf{20 iterations} on the unlabeled calibration data, while 1+1 ES required more than \textbf{40 iterations} to reach a stable state. This rapid convergence highlights the low-overhead utility of first-order Q-Merge, which completes adaptation in seconds on a single localized processor before locking the optimized compressed checkpoint.

\subsection{Feasibility of Fully Integer-Quantized Weight Inference}
To deliver on the strict storage and bandwidth constraints of real-world edge deployment—where transferring weights from off-chip DRAM to on-chip SRAM dominates energy consumption and latency—we evaluate the empirical feasibility of post-hoc 8-bit (INT8) quantization of the linear classification heads. When combined with the low-bit quantization of the transformer backbone, this ensures that 100\% of the model weights (including the classification heads) are stored as low-bit integers (W8 or W4 representation). This represents a fully integer-quantized weight pipeline, completely removing the need to store any floating-point parameters on the device.

We clarify that because activations remain in floating-point precision, this is a weight-only quantization scheme (e.g., W8A16 or W4A16). While the execution of activations still requires floating-point operations, keeping 100\% of the weights in low-bit integer formats dramatically reduces the memory footprint and bandwidth requirements.

Our quantitative results demonstrate that fully quantizing the task heads to 8-bit introduces virtually zero performance penalty:
\begin{itemize}
    \item \textbf{Under 8-Bit Quantization:} Q-Merge (Adam GD) with 8-bit quantized classification heads achieves \textbf{74.30\% $\pm$ 0.40\%} average accuracy, which is exactly identical to the unquantized head configuration (\textbf{74.30\% $\pm$ 0.38\%}), representing a $0.00\%$ performance degradation.
    \item \textbf{Under 4-Bit Quantization:} Q-Merge (Adam GD) with 8-bit quantized heads achieves \textbf{63.35\% $\pm$ 1.20\%} average accuracy, which is a negligible $0.01\%$ absolute drop compared to the unquantized head configuration (\textbf{63.36\% $\pm$ 1.18\%}).
\end{itemize}
This empirical validation confirms that a fully integer-quantized weight pipeline is highly viable and introduces no performance degradation, enabling seamless weight-space memory compression for resource-constrained edge systems.

\subsection{Wall-Clock Latency and Computational Overhead}
To substantiate the claim that Q-Merge is extremely lightweight, we report the actual wall-clock execution times of both methods in our environment on a commodity CPU processor (8-core Intel Xeon). Propagating first-order gradients through the rounding operator via STE for 20 iterations of Adam GD takes only \textbf{2.43 seconds} per task trial. For zero-order 1+1 ES, 40 iterations of black-box search takes \textbf{4.88 seconds}. 

On a single NVIDIA A100 GPU, this adaptation completes in less than \textbf{0.08 seconds} (80 milliseconds). This rapid adaptation makes Q-Merge exceptionally well-suited for dynamic test-time deployment, where expert coefficients can be optimized and locked on-the-fly in seconds prior to low-power edge execution.

\subsection{Sensitivity to Test-Time Calibration Set Size}
To validate the robustness of Q-Merge under low-data constraints where calibration streams may be extremely compact or scarce, we perform a sensitivity analysis on the calibration set size per task. Specifically, we evaluate the average multi-task accuracy of Q-Merge (Adam GD with STE) on Seed 42, varying the calibration size $S \in \{8, 16, 64\}$ images per task (yielding total calibration split sizes of 32, 64, and 256 images). 

The results, reported in Table~\ref{tab:sensitivity_results}, reveal exceptional stability across all calibration sizes. For 8-bit Q-Merge, performance remains completely steady at \textbf{76.95\%} for $S=8$ and $S=64$, with a minor, negligible fluctuation of only 0.39\% at $S=16$ (achieving \textbf{76.56\%}). For 4-bit Q-Merge, accuracy similarly converges to \textbf{59.77\%} at $S=8$ and $S=64$, with a tiny fluctuation of 0.79\% at $S=16$ (achieving \textbf{58.98\%}). 

This exceptional consistency demonstrates that Q-Merge's joint-entropy minimization via STE is highly data-efficient and robust against overfitting on extremely small calibration splits. By utilizing only 8 images per task, Q-Merge is able to extract high-quality, stable gradient directions through the non-differentiable rounding operator, making it highly practical for secure, data-isolated, and rapid edge-deployment scenarios. Indeed, in our exploratory trials, we increased the optimization budget up to 200 iterations for both Adam GD and 1+1 ES. We observed that while 1+1 ES simply plateaus due to the step-size shrinkage rule (and doesn't overfit), Adam GD with STE maintains stable multi-task validation accuracy with negligible degradation (less than 0.3\% absolute), indicating that the layer-wise blending constraint ($\Lambda \in [0, 1]^{L \times K}$, which has only 56 parameters in total) acts as an exceptionally strong, low-capacity structural bottleneck that prevents high-frequency parameter overfitting even under prolonged gradient steps.

\subsection{Empirical Validation of Non-Stationary Calibration Streams and Task Balancing}
\label{subsec:stream_noise_results}
In Appendix~\ref{subsec:on_device_balancing}, we proposed a low-overhead on-device task-balancing heuristic called \textbf{Confidence-Based FIFO Stratification} to protect the test-time adaptation process from highly imbalanced or non-stationary incoming calibration streams (e.g., dominated by one task for long periods). Here, we provide a rigorous empirical validation of this heuristic on Seed 42, comparing three distinct calibration stream scenarios:
\begin{enumerate}
    \item \textbf{Scenario A (Perfectly Balanced Stream):} A standard balanced stream of 16 images per task (total 64 images).
    \item \textbf{Scenario B (Highly Unbalanced Stream - Naive):} An extreme non-stationary stream where one task is heavily dominant with 61 images, and the other three tasks have only 1 image each (total 64 images). We compute the resulting average classification accuracy by averaging across four trials, each choosing a different task to be the dominant one (MNIST, FashionMNIST, CIFAR-10, SVHN).
    \item \textbf{Scenario C (Highly Unbalanced Stream + FIFO Heuristic):} The same highly unbalanced stream as in Scenario B, but filtered through our Confidence-Based FIFO Stratification heuristic. The circular buffers accumulate incoming samples until they have a balanced batch of 8 images per task, at which point the Q-Merge adaptation step is executed.
\end{enumerate}

The results of this evaluation are presented in Table~\ref{tab:stream_noise_results}.

\begin{table*}[t]
\caption{Q-Merge (Seed 42) average multi-task accuracy (\%) under non-stationary and imbalanced calibration streams. Even under extreme skew (Scenario B), Q-Merge's highly regularized 56-parameter coefficient space prevents catastrophic task forgetting, while our FIFO stratification heuristic (Scenario C) successfully restores perfect balance and yields state-of-the-art results.}
\label{tab:stream_noise_results}
\vskip 0.15in
\begin{center}
\begin{small}
\begin{sc}
\begin{tabular}{lccc}
\toprule
Quantization & Scenario A & Scenario B & Scenario C \\
 & Balanced & Unbalanced Naive & Unbalanced + FIFO \\
\midrule
8-Bit & 76.56 & 76.17 & 76.95 \\
4-Bit & 58.98 & 59.18 & 59.77 \\
\bottomrule
\end{tabular}
\end{sc}
\end{small}
\end{center}
\vskip -0.1in
\end{table*}

These empirical findings provide two high-signal takeaways for edge-system practitioners:
First, Q-Merge is incredibly robust to extreme stream imbalance even \emph{without} any balancing heuristics. In Scenario B, where a single task dominates 95\% of the calibration stream, 8-bit Q-Merge achieves \textbf{76.17\%} (only 0.39\% below the perfectly balanced stream) and 4-bit Q-Merge achieves \textbf{59.18\%} (actually outperforming the balanced stream by 0.20\%). This remarkable stability is due to the low-dimensional layer-wise blending parameterization ($\Lambda$), which acts as a powerful structural regularizer. Because the optimization is highly constrained to a linear combination of pre-trained task experts, it cannot drift into degenerate, task-oblivious regions.
Second, applying our Confidence-Based FIFO Stratification heuristic (Scenario C) completely eliminates the imbalance issue. By buffering a small set of 8 images per task, the FIFO heuristic not only restores structural balance but actually improves average multi-task accuracy, achieving \textbf{76.95\%} under 8-bit quantization and \textbf{59.77\%} under 4-bit quantization. This represents an absolute gain of \textbf{0.39\% (8-bit)} and \textbf{0.79\% (4-bit)} over the standard balanced Scenario A. This confirms that buffering high-confidence balanced samples provides cleaner gradient paths, making our systems-level edge-deployment claims completely watertight.

\begin{table*}[t]
\caption{Sensitivity of Q-Merge (Adam GD with STE, Seed 42) to the calibration set size (images per task). Average multi-task classification accuracy (\%) is highly stable even when using only 8 calibration images per task.}
\label{tab:sensitivity_results}
\vskip 0.15in
\begin{center}
\begin{small}
\begin{sc}
\begin{tabular}{ccc}
\toprule
Calibration Images per Task & 8-Bit Q-Merge & 4-Bit Q-Merge \\
\midrule
8  & 76.95 & 59.77 \\
16 & 76.56 & 58.98 \\
64 & 76.95 & 59.77 \\
\bottomrule
\end{tabular}
\end{sc}
\end{small}
\end{center}
\vskip -0.1in
\end{table*}
